\chapter{СИСТЕМИЙН ЗОХИОМЖ}
\begin{sloppy}

\section{Хэрэглэгчийн ба системийн шаардлагын шинжилгээ}

NomadAdventure системийн зохиомжийг боловсруулахдаа аялал жуулчлалын хэрэглэгчийн зан төлөв, байршлын системийн техникийн хязгаарлалт, AR орчны найдвартай ажиллагаа, соёлын контентыг тоглоомчилсон хэлбэрээр хүргэх шаардлагыг зэрэг авч үзэв. Энэхүү систем нь энгийн мэдээллийн сан үзүүлэх аппликейшн биш; харин хэрэглэгчийг бодит байршил дээр хөдөлж, AR орчинд интерактив байдлаар оролцож, reward collection хийх замаар танин мэдэхүйн туршлага олгох зорилготой.

\subsection{Функциональ шаардлага}

Систем дараах үндсэн үйл ажиллагааг дэмжих шаардлагатай.

\begin{itemize}
    \item Хэрэглэгчийн одоогийн байршлыг тодорхойлох, accuracy мэдээллийг шалгах.
    \item Аялал жуулчлалын байршлуудын жагсаалтыг backend-ээс татах болон шаардлагатай үед кешээс сэргээх.
    \item Хэрэглэгч тухайн байршлын радиуст орсон эсэхийг тооцож AR нээх боломжийг идэвхжүүлэх.
    \item AR орчныг эхлүүлж, виртуал авдрыг боломжит хавтгай дээр эсвэл fallback байрлалд дүрслэх.
    \item Авдар нээгдэхэд reward мэдээллийг үзүүлэх болон reward claim-ийг backend дээр баталгаажуулах.
    \item Хэрэглэгчийн цуглуулсан шагналуудыг collection дэлгэцээр харах.
\end{itemize}

\subsection{Функциональ бус шаардлага}

Системийн чанар, тогтвортой ажиллагааг дараах шаардлагаар тодорхойлов.

\begin{itemize}
    \item GPS савлалттай үед ч UX тогтвортой байх.
    \item Plane detection амжилтгүй үед AR туршлага бүрэн зогсохгүй fallback механизмаар үргэлжлэх.
    \item Давхар reward claim үүсэхгүй байх.
    \item Хэрэглэгчийн интерфэйс бага алхамтай, хөдөлгөөн дунд ашиглахад ойлгомжтой байх.
    \item Сүлжээ тасарсан үед сүүлийн татсан байршлын мэдээллээр үндсэн урсгал хязгаарлагдмал түвшинд ажиллах.
\end{itemize}

\subsection{Шаардлагын аналитик тайлбар}

Эдгээр шаардлагыг нэгтгэж үзвэл системийн амжилтыг зөвхөн ``AR ажиллаж байна уу'' гэдгээр бус, харин байршлын логик, UX-ийн ойлгомжтой байдал, өгөгдлийн бүрэн бүтэн байдал, fallback тэсвэржилт зэрэг нийлмэл хэмжээсээр үнэлэх шаардлагатай. Иймээс зохиомжийн шийдвэр бүрийг хэрэглэгчийн туршлага ба инженерийн эрсдэл хоёулантай нь уялдуулан гаргасан.

\section{Тоглоомын механик ба хэрэглэгчийн урсгал}

\subsection{Тоглоомын үндсэн цикл}

Системийн тоглоомын механик нь treasure hunt төрлийн энгийн боловч тодорхой ахицын логиктой. Хэрэглэгч эхлээд ойролцоох аяллын цэгийг илрүүлнэ, дараа нь тухайн байршлын радиуст орж AR горимыг нээнэ, AR орчинд авдрыг илрүүлж нээнэ, reward авсны дараа collection-д хадгална. Энэхүү цикл нь exploration, discovery, reward, cultural interpretation гэсэн дөрвөн түвшний туршлага үүсгэнэ.

\subsection{Reward collection ба соёлын өгүүлэмж}

Reward бүр нь Монголын ахуй, түүх, уламжлалтай уялдсан тайлбартай байна. Өөрөөр хэлбэл шагнал нь зөвхөн дүрслэлийн объект бус, тухайн байршлын танин мэдэхүйн өгөгдлийг агуулсан контентын нэгж юм. Collection дэлгэцэд reward-уудыг авсан, аваагүй төлөвөөр ялган харагдуулснаар ахицын мэдрэмж төрүүлнэ. Нэг байршлаас нэг reward зөвхөн нэг удаа авах дүрэм нь exploration-г шинэ байршил руу чиглүүлнэ.

\subsection{Хэрэглэгчийн интерфэйсийн урсгал}

Системийн интерфэйсийг дараах үндсэн дэлгэцүүдэд төвлөрүүлэв: газрын зургийн дэлгэц, AR дэлгэц, reward panel, collection дэлгэц. Газрын зургийн дэлгэц дээр ойролцоох аяллын цэг, радиуст орсон эсэхийн төлөв харагдана. AR дэлгэц нь камер ба AR session-ийг ажиллуулж авдрыг дүрслэнэ. Reward panel нь тухайн шагналын нэр, тайлбар, claim төлөвийг харуулна. Collection дэлгэц нь өмнө нь авсан шагналуудыг нэгтгэнэ.

Текстэн зураглалын тайлбар: хэрэглэгч аппликейшн нээхэд guest session үүснэ. Дараа нь аяллын цэгүүд ачаалагдаж газрын зурагт дүрслэгдэнэ. Хэрэглэгч радиуст ормогц AR товч идэвхжинэ. AR scene-д авдрыг нээсний дараа reward panel гарч, claim амжилттай бол collection state шинэчлэгдэн газрын зураг эсвэл collection дэлгэц рүү буцна.

\subsection{UX-ийн зохиомжийн үндэслэл}

Аялал жуулчлалын орчинд хэрэглэгч байнга алхаж, орчноо ажиглаж, бусадтай харилцаж байдаг тул урсгал богино байх ёстой. Иймээс энэ систем төвөгтэй даалгавар, олон шатат мини тоглоом, урт тайлбараас зайлсхийж, ``хүрэх--нээх--авах--унших'' гэсэн ойлгомжтой UX-ийг баримталсан. Энэ нь flow-ийг алдагдуулахгүй, мөн GPS ба AR-ийн техникийн эрсдэлийг хэрэглэгчийн ойлголтын хувьд төвөгтэй болгохгүй байх давуу талтай.

\section{Архитектурын шийдэл ба AR идэвхжүүлэх механизм}

\subsection{Ерөнхий архитектур}

Систем нь Unity төвтэй мобайл клиент, Node.js/Express backend, PostgreSQL/Prisma өгөгдлийн сан гэсэн гурван үндсэн хэсэгтэй. Клиент тал байршлын шинэчлэлт, газрын зураг, AR scene, collection UI-г хариуцна. Backend нь аяллын цэгийн мета мэдээлэл, reward claim, collection retrieval, guest session-ийг удирдана. Өгөгдлийн сан нь хэрэглэгчийн төлөв, аяллын цэг, reward, claim түүхийг хадгална.

\subsection{Макро ба микро түвшний логикийн салгалт}

Зохиомжийн гол инженерийн шийдвэр нь байршлын идэвхжүүлэлтийг макро түвшинд, AR байршуулалтыг микро түвшинд салгаж үзсэнд оршино. Макро түвшинд GPS координат, Haversine зай, smoothing, hysteresis ашиглан хэрэглэгч радиуст орсон эсэхийг тодорхойлно. Микро түвшинд AR session эхэлсний дараа camera tracking, plane detection, anchor placement ашиглан авдрыг орчны боломжит хавтгай дээр байрлуулна.

Энэ салгалтын гол давуу тал нь GPS-ийн хэлбэлзэл виртуал объектын байрлалд шууд нөлөөлөхгүй байх явдал юм. Хэрэв бодит координатыг шууд 3D объектод шилжүүлбэл хэрэглэгч авдрыг хайж төөрөх, эсвэл объект харагдахгүй байх эрсдэл нэмэгдэнэ. Динамик placement нь хэрэглэгчид илүү ойлгомжтой, prototype түвшинд найдвартай UX өгнө.

\subsection{AR идэвхжүүлэх логик}

Байршлын шинэчлэлт ирэх бүрд систем ойролцоох цэгүүдийн зайг шалгана. Радиусын `enter` босгыг давсан үед AR товч идэвхжиж, `exit` босгыг давмагц идэвхгүй болно. Хэрэв хэрэглэгч өмнө нь тухайн reward-ийг claim хийсэн бол AR товчийг мэдээллийн горимд шилжүүлж, reward дахин олгохгүй. Ийм логик нь өгөгдлийн түвшний давхардал хамгаалалттай нийцэж ажиллана.

\subsection{Fallback ба resilience дизайн}

Систем дараах fallback механизмыг зохиомжийн түвшинд тусгав.

\begin{itemize}
    \item Аяллын цэгийн өгөгдөл татагдахгүй бол сүүлийн кешлэгдсэн хувилбарыг ашиглах.
    \item Plane detection хангалтгүй бол авдрыг хэрэглэгчийн өмнө fallback байрлалаар дүрслэх.
    \item Reward зураг эсвэл prefab олдохгүй бол placeholder дүрслэл үзүүлэх.
    \item Сүлжээ тасалдсан үед reward claim-ийг шууд хийхгүй, харин хэрэглэгчид тодорхой тайлбар өгөх.
\end{itemize}

Энэхүү fallback дизайн нь UX-ийг таслахгүй байлгахын зэрэгцээ хэрэглэгчид системийн төлөвийг ойлгомжтой тайлбарлах боломжийг олгоно.

\section{Өгөгдлийн загвар ба бизнес дүрэм}

\subsection{Entity-үүдийн үүрэг}

\begin{table}[H]
\centering
\caption{Системийн үндсэн entity-үүд ба үүрэг}
\begin{tabular}{|p{3cm}|p{9cm}|}
\hline
\textbf{Entity} & \textbf{Үүрэг} \\ \hline
User & Guest session, display name, төхөөрөмжийн танигч, collection-ийн эзэн \\ \hline
TourismSpot & Аяллын цэгийн нэр, координат, радиус, reward холбоос \\ \hline
Reward & Шагналын нэр, тайлбар, дүрслэлийн мета мэдээлэл \\ \hline
UserReward & Тухайн хэрэглэгч тухайн цэгийн reward-ийг хэдийд claim хийснийг хадгалах холбоос хүснэгт \\ \hline
\end{tabular}
\end{table}

\subsection{Business rule-үүд}

Системийн гол бизнес дүрэм нь дараах байдалтай.

\begin{itemize}
    \item Нэг хэрэглэгч нэг байршлын reward-ийг нэг л удаа claim хийнэ.
    \item Reward claim амжилттай болсон үед collection retrieval-ийн дараагийн үр дүнд туссан байна.
    \item Reward claim өмнө нь бүртгэгдсэн бол сервер idempotent хариу буцаана.
    \item Хэрэглэгчийн байршил радиуст орсон ч AR session эхлэх техникийн нөхцөл хангагдаагүй бол мэдээллийн fallback горим ажиллана.
\end{itemize}

\subsection{API гэрээ}

Клиент ба серверийн хоорондох интерфейсийг дараах endpoint-ууд тогтооно.

\begin{itemize}
    \item \texttt{POST /auth/guest-login}: guest session үүсгэх буюу сэргээх
    \item \texttt{GET /spots}: идэвхтэй аяллын цэгүүдийг reward мэдээлэлтэй нь авах
    \item \texttt{POST /spots/:id/claim}: reward claim баталгаажуулах
    \item \texttt{GET /users/:id/rewards}: хэрэглэгчийн collection-ийг авах
\end{itemize}

Эдгээр гэрээ нь дипломын ажлын хэрэгжилтийн хэсэгт тайлбарлагдах класс, модуль, өгөгдлийн схемтэй шууд уялдана.

\section{Зохиомжийн аналитик дүгнэлт}

Энэхүү зохиомжийн бүлгийн дүн шинжилгээнээс харахад системийн амжилт нь зөвхөн UI-ийн гоо зүй, эсвэл AR-ийн эффектээс бус, харин байршлын логик, fallback стратеги, өгөгдлийн бүрэн бүтэн байдал, тоглоомжуулалтын механикийн уялдаанаас шалтгаална. Unity төвтэй архитектур сонгосон нь AR scene ба орон зайн харилцан үйлдлийг нэг орчинд төвлөрүүлж, клиент талын логикийг илүү нягт нэгтгэх давуу талтай.

Мөн ``радиуст суурилсан идэвхжилт + динамик байршуулалт + unique claim хамгаалалт'' гэсэн хосолсон шийдэл нь тухайн системийн гол инженерийн санаа болж байна. Энэ шийдэл нь онолын үндэслэлтэй, прототипийн түвшинд хэрэгжих боломжтой, цаашид өргөтгөхөд тохиромжтой бүтэц бүрдүүлж байгаа тул дараагийн бүлэгт эдгээрийг бодит хэрэгжилт ба туршилтын нотолгоогоор тайлбарлах боломжтой болж байна.

\end{sloppy}
